\section{Discussion and future work}
\label{sec:conclusion}
\seckiny

How to create agents that can learn to decompose their behaviour into meaningful primitives and then reuse them to more efficiently acquire new behaviours is a long standing research question. 
%
The solution to this question may be an important stepping stone towards agents with general intelligence and competence. %
This paper introduced FeUdal Networks, a novel architecture that formulates sub-goals as directions in latent state space, which, if followed, translate into a meaningful behavioural primitives. 
%
FuN clearly separates the module that discovers and sets sub-goals from the module that generates the behaviour through primitive actions. This creates a natural hierarchy that is stable and allows both modules to learn in complementary ways. 
%
Our experiments clearly demonstrate that this makes long-term credit assignment and memorisation more tractable. 
%
This also opens many avenues for further research, for instance: deeper hierarchies can be constructed by setting goals at multiple time scales, scaling agents to truly large environments with sparse rewards and partial observability.
%
The modular structure of FuN is also lends itself to transfer and multitask learning -- learnt behavioural primitives can be re-used to acquire new complex skills, or alternatively the transitional policies of the Manager can be transferred to agents with different embodiment.  

\section{Acknowledgements}
We thank Alex Graves, Daan Wierstra, Olivier Pietquin, Oriol Vinyals, Joseph Modayil and Vlad Mnih for many helpful
discussions, suggestions and comments on the paper.